\section{Ograniczenia badania}
Przeprowadzone badania posiadały pewne ograniczenia wynikające z przyjętych założeń oraz limitów sprzętowych. Z tego powodu badanie obejmowało ograniczoną liczbę modeli złożonych z 7 lub 8 miliardów parametrów. Otrzymane rezultaty, czyli skuteczność ataków, obrony, użyteczności oraz opóźnienia odzwierciedlają zatem system korzystający z takich modeli, działający na pojedynczym węźle obliczeniowym. Zakres badania ograniczono również do bezpośrednich ataków promptowych. Nie analizowano ataków pośrednich, w których system przyjmuje dane wejściowe ze źródeł zewnętrznych. Analiza i wnioski dotyczą zatem wyłącznie odporności modeli na dane wejściowe przekazywane w sposób bezpośredni. W pracy nie uwzględniono wszystkich kategorii mechanizmów obronnych. Zakres ograniczono do metod stosowanych na etapie wdrożenia modelu i pominięto podejścia wymagające trenowania lub dostrajania \gls{LLM}.

Badanie ograniczono do poleceń pochodzących ze zbioru JBB-Behaviours~\cite{chao2024jailbreakbenchopenrobustnessbenchmark}, ze względu na fakt, że ilość sposobów ominięcia zabezpieczeń przy użyciu promptu wejściowego jest w teorii nieskończona. Analizę skuteczności obrony przeprowadzono przy użyciu jedynie trzech odmiennych metodologicznie od siebie ataków. Nie ukazuje to jednak pełnego spektrum możliwych zagrożeń. Atak \gls{DAN} został przeprowadzony przy użyciu konkretnego, pojedynczego szablonu. W rzeczywistości istnieją różne warianty jego postaci. Zatem wyniki odporności na to zagrożenie odzwierciedlają zachowanie modeli wobec konkretnego szablonu \gls{DAN}, a nie całej klasy ataków opartych na odgrywaniu ról i wprowadzaniu fikcyjnego kontekstu. Należy także zaznaczyć, że konwersacje z modelem docelowym miały charakter jednoturowy. Model otrzymywał szkodliwy prompt wejściowy, na który generował odpowiedź. Nie rozpatrzono skuteczności ataków wieloturowych takich jak Crescendo~\cite{russinovich2025greatwritearticlethat}. Ponadto, wykorzystane złośliwe polecenia były jedynie w języku angielskim. Brak różnorodności w języku promptów stanowi lukę badawczą, która wymagałaby sprawdzenia w przyszłych eksperymentach. Doświadczenia przeprowadzono również na określonym limicie długości odpowiedzi i nie weryfikowano wpływu jego zwiększenia na skuteczność rozpatrywanych metod.

Otrzymana skuteczność poszczególnych metod ataku mogła wynikać z ograniczeń zasobów obliczeniowych. W przypadku techniki \gls{GCG} postanowiono na ograniczenie liczby kroków tej metody, co mogło wpłynąć na finalną jakość otrzymanych sufiksów. Nie skorzystano również z tego powodu z optymalizacji wejściowej pożądanej odpowiedzi na bazie AdvPrefix~\cite{zhu2025advprefixobjectivenuancedllm}, co mogło zapewnić większą skuteczność ataku \gls{GCG} i polepszyć ją w zależności od rozpatrywanej rodziny modeli. Wykorzystano zatem prefiksy dołączone do szkodliwych instrukcji w zbiorze promptów wejściowych. Przerostki te mogły być nadmiernie ograniczone, co zauważono w badaniu Zhu i in.~\cite{zhu2025advprefixobjectivenuancedllm}.

Skuteczność metody \gls{PAIR} zależała od doboru modelu atakującego oraz modelu oceniającego wygenerowaną odpowiedź, więc otrzymane wyniki były zależne od ich ogólnej zdolności. Użycie modeli o większej liczbie parametrów i brak konieczności kwantyzacji wag mógłby wpłynąć na poprawę skuteczności tego ataku. Zmniejszono limit iteracji, by zoptymalizować czas trwania eksperymentu. Znaczne zwiększenie tego limitu mogło by skutkować znalezieniem większej liczby obejść wdrożonych zabezpieczeń modelu.

Kluczowym w fazie inferencji było zastosowanie deterministycznego zachłannego dekodowania oraz wyzerowanie parametru temperatury modelu. Choć decyzja ta była konieczna do zapewnienia powtarzalność pomiarów, ograniczyła wariancję odpowiedzi. Istnieje możliwość, że dla innej konfiguracji z losowym próbkowaniem wraz z wyższą wartością parametru temperatury otrzymano by odmienne rezultaty.
Wpływ na finalne wyniki mogło mieć również zastosowanie \gls{LLM} do oceny generownych odpowiedzi przez model docelowy. Zrezygnowano z manualnej ewaluacji na rzecz klasyfikacji automatycznej. Podejście to zagwarantowało powtarzalność i skalowalność, lecz model weryfikujący mógł posiadać własne uprzedzenia, co w niewielkim stopniu mogło wpłynąć na finalne rezultaty wskaźnika \gls{ASR}.

W przypadku metody opartej na Llama-Guard zastosowano ocenę pełnej odpowiedzi modelu docelowego przed jej zwróceniem, co mogło ograniczyć responsywność badanego systemu. Rozważyć można by wariant oparty na dzieleniu generowanej odpowiedzi i klasyfikacji na bieżąco zwracanych części tekstu.